%!TEX root = ../template.tex
%%%%%%%%%%%%%%%%%%%%%%%%%%%%%%%%%%%%%%%%%%%%%%%%%%%%%%%%%%%%%%%%%%%%
%% abstract-en.tex
%% NOVA thesis document file
%%
%% Abstract in English([^%]*)
%%%%%%%%%%%%%%%%%%%%%%%%%%%%%%%%%%%%%%%%%%%%%%%%%%%%%%%%%%%%%%%%%%%%

\typeout{NT FILE abstract-en.tex}%

\par The Web3 movement is currently pushing for a more decentralized (and democratic) architecture for the Internet. As such, decentralized systems are gaining popularity once more. A particular case of peer-to-peer usage appears in Swarm Computing, where highly heterogeneous devices can cooperate directly and in a mostly autonomous fashion, demands for highly flexible and robust decentralized solutions, including distributed and decentralized membership, as well as communication protocols\todo{"Swarm computing Airdropped" what does this mean? Should I just remove the reference?}.
\par The peer-to-peer literature is full of different proposals for such protocols. However, the majority assumes a static configuration, meaning that the parameters that govern the operation of the system are defined at bootstrap time, based on expectations on the operational conditions of the system throughout its entire lifetime. These remain unchanged even if these conditions change significantly, making the system potentially inefficient or even unavailable. The autonomic computing paradigm has been proposed many years ago as a way to achieve a generic architecture to build systems that monitor themselves and their environment to make decisions on how to reconfigure themselves, and apply those reconfiguration actions. Solutions focused on autonomic computing are usually fully decentralized, or fully centralized, leaving space for the study of hybrid solutions.
\par In this work we design a centralized (or partially centralized) component that can make management decisions for a large-scale decentralized system, which we named Overlord. We explore the idea of autonomic computing to govern peer-to-peer protocol parameters. To this end, we create a design for the architecture, implement it using the Babel-Swarm framework, and evaluate its performance against static configuration alternatives in a message dissemination case study.

\keywords{
  Decentralized Systems \and
  Autonomic Computing \and
  Peer-To-Peer \and
  Self-Adaptive Systems
}
